\section[Cote 74 : complexes cubiques généralisés, la réduction à un carré latin et à ses revêtements doubles]{Cote 74 : complexes cubiques généralisés, la réduction à un carré latin et à ses revêtements doubles\protect\verifie{Folder74}}\label{sec:74}

La cote 74, \emph{Complexes cubiques : notes manuscrites (s.d.)}, cent trente-cinq pages que l'inventaire date \q{[à partir de 1976]}, tourne autour du graphe des vingt-sept droites d'une surface cubique. La liasse qui commence page~92 porte le titre \q{Complexes cubiques généralisés (complexes \emph{triadiques} ou \emph{trilatères} ?)}. Elle pose deux axiomes sur un graphe $S$ : toute arête est contenue dans un unique triangle (axiome~1), et tout sommet hors d'un triangle est lié à un sommet de ce triangle (axiome~2). Elle fixe un triangle $t_0 = \{a, b, c\}$ et décompose le graphe autour de lui (\lot{74}{5}{92}). Les pages~93 à~96 établissent les propositions~1 à~4 et la \q{Scholie} : le graphe se résume à trois ensembles $\widetilde E_\alpha$ munis d'involutions sans point fixe et à une partie $\widetilde T$ de leur produit, soumise aux conditions~(*) et~(**). Une troisième condition, notée~(***), est écrite en travers de la marge gauche, encadrée, tournée de 90°, et rattachée par une accolade à~(*) et~(**). La transcription en lit le début, \q{si $(s_\alpha)_{\alpha\in t_0}\in\prod\widetilde{E}_\alpha$ et $(t_\alpha)_{\alpha\in t_0}\in\widetilde{T}$ tels que $s_\beta=\varphi_\beta(t_\alpha)$ \ldots{} $\forall\,\alpha,\beta$, $\alpha\neq\beta$}, et la fin est illisible (\lot{74}{5}{96}). La page~97 affirme que des données vérifiant~(*) donnent un graphe qui vérifie l'axiome~1 et l'axiome~2 pour $t_0$ (\lot{74}{5}{97}). Les pages~98 à~100 traitent l'axiome~2 pour les autres triangles : proposition~5 pour les triangles verticaux, axiome~3 (les $\widetilde E_\alpha$ sont non vides), proposition~6, barrée mais lisible, pour les triangles supérieurs (\lot{74}{5}{100}). La page~101 reprend la liste des données avec, en~($\gamma$), \q{l'axiome (Ka) plus haut}. Elle passe aux quotients $\Gamma = \widetilde\Gamma/\sigma$, $E_\alpha = \widetilde E_\alpha/\sigma_\alpha$ (\lot{74}{6}{101}). La page~102 énonce l'axiome $\beta_\Gamma$, \q{$\Gamma\to E_\alpha\times E_\beta$ est \emph{bijectif}}, et partage la classification en deux questions, I) les systèmes de trois quotients vérifiant $\beta_\Gamma$, II) leurs revêtements d'ordre~2, puis ceux qui satisfont (Ka) (\lot{74}{6}{102}, \lot{74}{6}{103}). La page~116 étudie, sur la relation $xyz = 1$ d'un groupe $G$, les applications $\varphi_a$ et conclut que tout élément est d'ordre~2 (\lot{74}{6}{116}).

\noindent\emph{Conventions.} Un graphe est sans boucle ni arête multiple ; un triangle est un ensemble de trois sommets deux à deux liés. On indexe $t_0$ par $\{0, 1, 2\}$, $t_0 = \{t_0(\alpha)\}$. Pour $\alpha \in t_0$, $\widetilde E_\alpha$ est l'ensemble des sommets hors de $t_0$ liés à $t_0(\alpha)$. Un triangle supérieur est un triplet $x = (x_\alpha)$ avec $x_\alpha \in \widetilde E_\alpha$ et les $x_\alpha$ deux à deux liés ; leur ensemble est $\widetilde\Gamma$ (le $\widetilde T$ de la page~96). Pour des données abstraites $(\widetilde E_\alpha, \sigma_\alpha, \widetilde T)$, deux points $s \in \widetilde E_\alpha$ et $u \in \widetilde E_\beta$, $\alpha \neq \beta$, sont \emph{liés} s'il existe $\tau \in \widetilde T$ avec $\tau_\alpha = s$ et $\tau_\beta = u$. Le graphe des données a pour sommets $t_0 \amalg \coprod_\alpha \widetilde E_\alpha$ et pour arêtes celles de 2°)~$\alpha$) à~$\delta$) de la page~96 : deux sommets de $t_0$, $t_0(\alpha)$ et les points de $\widetilde E_\alpha$, $s$ et $\sigma_\alpha s$, deux points liés. Les conditions sur les données sont :
\begin{itemize}
\item[(*)] deux éléments de $\widetilde T$ qui ont deux coordonnées communes sont égaux ;
\item[($\beta_{\widetilde\Gamma}$)] pour $\alpha \neq \beta$, $s \in \widetilde E_\alpha$ et $u \in \widetilde E_\beta$, un et un seul des deux points $s$, $\sigma_\alpha s$ est lié à $u$ (c'est la forme de la page~101, qui contient~(**)) ;
\item[(Ka)] pour $\{\alpha, \beta, \gamma\} = t_0$ et $\tau, \rho_1, \rho_2 \in \widetilde T$ tels que $\tau_\alpha = (\rho_1)_\alpha$, $\tau_\beta = (\rho_2)_\beta$ et $(\rho_1)_\gamma = (\rho_2)_\gamma$, on a $\tau = \rho_1$.
\end{itemize}
La condition (Ka) ainsi écrite est une restitution. Elle porte sur une famille $(s_\alpha)$ et trois triplets de $\widetilde T$, chacun contenant deux des $s_\alpha$, comme le début lisible de la marge. On note (Ka$'$) la forme qui garde la phrase de la marge et la complète par la conclusion $(s_\alpha) \in \widetilde T$ : si $(s_\alpha) \in \prod \widetilde E_\alpha$ et $(t_\alpha) \in \widetilde T^{t_0}$ vérifient $(t_\alpha)_\beta = s_\beta$ pour $\alpha \neq \beta$, alors $(s_\alpha) \in \widetilde T$.

\begin{enonce}[Pages 92 à 99, autour d'un triangle\piste{74-gq2t-latin-square-reduction}]
Soit $S$ un graphe vérifiant les axiomes~1 et~2, et $t_0$ un triangle de $S$.
\begin{enumerate}
\item (Page~92, (1)) Tout sommet hors de $t_0$ est lié à un et un seul sommet de $t_0$ ; les $\widetilde E_\alpha$ sont disjoints et recouvrent $S \smallsetminus t_0$.
\item (Prop.~1) Pour $s \in \widetilde E_\alpha$, le troisième sommet $\sigma_\alpha s$ du triangle de l'arête $\{t_0(\alpha), s\}$ est dans $\widetilde E_\alpha$, et c'est le seul point lié à la fois à $t_0(\alpha)$ et à $s$. L'application $\sigma_\alpha$ est une involution sans point fixe de $\widetilde E_\alpha$.
\item (Prop.~2 et~5) Pour $\beta \neq \alpha$ et $u \in \widetilde E_\beta$, $u$ est lié à un et un seul des points $s$, $\sigma_\alpha s$.
\item (Prop.~3 et~4) Si $s \in \widetilde E_\alpha$ et $u \in \widetilde E_\beta$, $\alpha \neq \beta$, sont liés, le troisième sommet du triangle de l'arête $\{s, u\}$ est dans $\widetilde E_\gamma$, $\gamma$ le troisième indice ; il existe donc un triangle supérieur $x$ avec $x_\alpha = s$, $x_\beta = u$, et il est unique.
\item (Page~98, corollaires) Si $s$ et $u$ sont liés, $\sigma_\alpha s$ et $\sigma_\beta u$ le sont ; $\widetilde\Gamma$ est stable par $\sigma = \prod_\alpha \sigma_\alpha$, qui y est sans point fixe.
\end{enumerate}
\end{enonce}

\begin{enonce}[Pages 99, 101 et 102, le carré latin et les revêtements doubles\piste{74-gq2t-latin-square-reduction}]
Sous les mêmes hypothèses, soient $E_\alpha = \widetilde E_\alpha/\sigma_\alpha$, $\Gamma = \widetilde\Gamma/\sigma$ et $\varphi_\alpha : \Gamma \to E_\alpha$ l'application induite par $x \mapsto x_\alpha$.
\begin{enumerate}
\item (Axiome $\beta_\Gamma$) Pour $\alpha \neq \beta$, l'application $(\varphi_\alpha, \varphi_\beta) : \Gamma \to E_\alpha \times E_\beta$ est bijective : $\Gamma$ est un carré latin sur les trois $E_\alpha$.
\item (Page~101) Pour tout $\alpha$, l'application $x \mapsto (\bar x, x_\alpha)$ est une bijection de $\widetilde\Gamma$ sur le produit fibré $\Gamma \times_{E_\alpha} \widetilde E_\alpha$ : $\widetilde\Gamma$ est l'image inverse du revêtement double $\widetilde E_\alpha \to E_\alpha$ par $\varphi_\alpha$.
\item (Page~102) Les composés $\psi_{\beta\alpha} : \Gamma \times_{E_\alpha} \widetilde E_\alpha \to \Gamma \times_{E_\beta} \widetilde E_\beta$ de ces bijections sont des isomorphismes au-dessus de $\Gamma$, et forment un système transitif : $\psi_{\gamma\beta} \circ \psi_{\beta\alpha} = \psi_{\gamma\alpha}$, $\psi_{\alpha\alpha} = \mathrm{id}$.
\item (Pages~99 et~104) $|\widetilde E_\alpha| = 2\,|E_\alpha|$, $|\widetilde\Gamma| = 2\,|\Gamma|$, et $|\Gamma| = |E_\alpha|\,|E_\beta|$ pour $\alpha \neq \beta$. Si $\widetilde E_\alpha$ est non vide, $E_\beta$ et $E_\gamma$ sont en bijection.
\end{enumerate}
\end{enonce}

\begin{enonce}[Pages 96, 97 et 100, la Scholie et sa réciproque\piste{74-gq2t-latin-square-reduction}]
\begin{enumerate}
\item (Scholie) Le graphe $S$ est isomorphe au graphe de ses données $(\widetilde E_\alpha, \sigma_\alpha, \widetilde\Gamma)$, et ces données vérifient~(*), $\beta_{\widetilde\Gamma}$ et~(Ka).
\item (Réciproque) Pour des données vérifiant $\beta_{\widetilde\Gamma}$ et~(Ka), le graphe des données vérifie les axiomes~1 et~2, et $t_0$ en est un triangle. Ses triangles supérieurs sont les éléments de $\widetilde T$ : trois points deux à deux liés hors de $t_0$ sont les trois coordonnées d'un $\tau \in \widetilde T$.
\item Pour des données vérifiant~(*) et $\beta_{\widetilde\Gamma}$, le graphe vérifie l'axiome~1 si et seulement si~(Ka) est vraie.
\item (Ka) entraîne~(*), et (Ka) équivaut à la conjonction de~(*) et de~(Ka$'$).
\end{enumerate}
\end{enonce}

\begin{enonce}[Pages 101 à 103, problèmes I et II\piste{74-gq2t-latin-square-reduction}]
Soient $\Gamma$ un ensemble, $\varphi_\alpha : \Gamma \to E_\alpha$ ($\alpha \in t_0$) trois applications vérifiant $\beta_\Gamma$, $p_\alpha : \widetilde E_\alpha \to E_\alpha$ trois revêtements doubles, c'est-à-dire des applications dont les fibres sont les orbites d'une involution sans point fixe $\sigma_\alpha$, et $\widetilde\Gamma$ un ensemble muni de $\pi : \widetilde\Gamma \to \Gamma$ et de $\psi_\alpha : \widetilde\Gamma \to \widetilde E_\alpha$ tels que $p_\alpha \psi_\alpha = \varphi_\alpha \pi$ et que $(\pi, \psi_\alpha) : \widetilde\Gamma \to \Gamma \times_{E_\alpha} \widetilde E_\alpha$ soit bijective pour tout $\alpha$. Alors les données $(\widetilde E_\alpha, \sigma_\alpha, \widetilde T)$, $\widetilde T$ l'image de $(\psi_\alpha)_\alpha$, vérifient~(*) et $\beta_{\widetilde\Gamma}$. Si elles vérifient~(Ka), leur graphe vérifie les axiomes~1 et~2 ; il vérifie l'axiome~1 si et seulement si elles vérifient~(Ka). Un graphe muni d'un triangle $t_0$ fournit de telles données par l'énoncé précédent.
\end{enonce}

\begin{enonce}[Marge de la page 96, la condition (Ka) est nécessaire]
Soient $E_\alpha = \mathbb Z/2$, $\widetilde E_\alpha = \mathbb Z/2 \times \mathbb Z/2$, $\sigma_\alpha(u, e) = (u, e + 1)$, et $\widetilde T$ l'ensemble des triplets $((u, e), (v, e), (u + v, e))$. Ce sont les données du carré latin d'ordre~2, table de $\mathbb Z/2$, avec trois revêtements triviaux. Elles vérifient~(*) et $\beta_{\widetilde\Gamma}$, et ne vérifient pas~(Ka). Leur graphe ne vérifie pas l'axiome~1 : l'arête qui joint $(0,0) \in \widetilde E_a$ à $(0,0) \in \widetilde E_b$ est contenue dans deux triangles, fermés par $(0,0)$ et par $(1,0) \in \widetilde E_c$.
\end{enonce}

\begin{enonce}[Page 116\piste{74-gq2t-latin-square-reduction}]
Soit $G$ un groupe. Les applications $\varphi_a : (x, y, z) \mapsto (x, a^{-1}z^{-1}, a^{-1}y^{-1})$ préservent la relation $xyz = 1$ pour tout $a \in G$ si et seulement si $g^2 = 1$ pour tout $g \in G$. Le groupe $G$ est alors commutatif.
\end{enonce}

\begin{proof}[Démonstration des pages 92 à 99]
(1) Soit $s$ hors de $t_0$. L'axiome~2 pour $t_0$ donne un sommet de $t_0$ lié à $s$. S'il y en avait deux, $t_0(\alpha)$ et $t_0(\beta)$, l'arête $\{t_0(\alpha), t_0(\beta)\}$ serait contenue dans le triangle $t_0$ et dans le triangle qu'elle forme avec $s$ ; l'axiome~1 donnerait $s = t_0(\gamma)$, contre l'hypothèse.

(2) Soit $s \in \widetilde E_\alpha$ et $z$ le troisième sommet du triangle de $\{t_0(\alpha), s\}$. Si $z$ était un sommet $t_0(\beta)$ de $t_0$, on aurait $\beta \neq \alpha$ et $s$ serait lié à $t_0(\alpha)$ et à $t_0(\beta)$, contrairement à~(1) ; donc $z \in \widetilde E_\alpha$. Tout point lié à $t_0(\alpha)$ et à $s$ est égal à $z$ par l'axiome~1. Le point $s$ est lié à $t_0(\alpha)$ et à $\sigma_\alpha s$, donc $\sigma_\alpha \sigma_\alpha s = s$, et $\sigma_\alpha s \neq s$ puisque $s$ et $\sigma_\alpha s$ sont liés.

(3) Si $u$ était lié à $s$ et à $\sigma_\alpha s$, l'arête $\{s, \sigma_\alpha s\}$ serait contenue dans le triangle qu'elle forme avec $t_0(\alpha)$ et dans celui qu'elle forme avec $u$, d'où $u = t_0(\alpha)$, absurde (proposition~2 ; l'hypothèse $\beta \neq \alpha$ n'y sert pas). L'axiome~2 pour le triangle vertical $\{t_0(\alpha), s, \sigma_\alpha s\}$ lie $u$ à l'un de ses trois sommets ; ce n'est pas $t_0(\alpha)$ par~(1), puisque $u$ est lié à $t_0(\beta)$ ; c'est donc $s$ ou $\sigma_\alpha s$ (proposition~5).

(4) Soit $z$ le troisième sommet du triangle de $\{s, u\}$. Si $z$ était un sommet $t_0(\delta)$, (1) donnerait $\delta = \alpha$ et $\delta = \beta$. Par~(1), $z$ est dans un $\widetilde E_\delta$. Si $\delta = \alpha$, $z$ est lié à $t_0(\alpha)$ et à $s$, donc $z = \sigma_\alpha s$ par~(2), et $u$ est lié à $s$ et à $\sigma_\alpha s$, contrairement à~(3) ; de même $\delta \neq \beta$. Donc $\delta = \gamma$, et $(s, u, z)$, rangé selon les indices, est un triangle supérieur. Deux triangles supérieurs $x$, $y$ avec $x_\alpha = y_\alpha$ et $x_\beta = y_\beta$ ont $x_\gamma$ et $y_\gamma$ liés à $x_\alpha$ et à $x_\beta$, donc égaux par l'axiome~1.

(5) Si $s$ et $u$ sont liés, $\sigma_\beta u$ n'est pas lié à $s$ par~(3) appliqué au triangle vertical de $u$ ; il est donc lié à $\sigma_\alpha s$ par~(3) appliqué au triangle vertical de $s$. Les coordonnées de $\sigma x$ sont donc deux à deux liées, et $\sigma x \neq x$ puisque $\sigma_\alpha$ est sans point fixe.
\end{proof}

\begin{lemme}\label{lem:74-revetement}
Soit $X$ un ensemble muni d'une involution sans point fixe $f$. Le choix d'un représentant dans chaque orbite donne une bijection $X \simeq X/f \times \{0, 1\}$ ; en particulier $|X| = 2\,|X/f|$.
\end{lemme}

\begin{proof}
Soit $r(q)$ le représentant de l'orbite $q$. On envoie $x$ sur $(\bar x, 0)$ si $x = r(\bar x)$ et sur $(\bar x, 1)$ sinon, et $(q, 0)$ sur $r(q)$, $(q, 1)$ sur $f(r(q))$. Comme l'orbite de $x$ est $\{r(\bar x), f(r(\bar x))\}$ et que $f$ est sans point fixe, ces deux applications sont réciproques.
\end{proof}

\begin{lemme}\label{lem:74-ligne}
Soient $f : \Gamma \to A$, $g : \Gamma \to B$, $h : \Gamma \to D$ telles que $(f, g)$ et $(f, h)$ soient bijectives. Pour tout $a \in A$, l'application $b \mapsto h\bigl((f, g)^{-1}(a, b)\bigr)$ est une bijection de $B$ sur $D$.
\end{lemme}

\begin{proof}
Si deux valeurs $b_1$, $b_2$ ont même image, les éléments $(f, g)^{-1}(a, b_i)$ ont même $f$ et même $h$ ; ils sont égaux par l'injectivité de $(f, h)$, d'où $b_1 = b_2$. Pour $d \in D$, soit $x = (f, h)^{-1}(a, d)$ ; alors $(f, g)^{-1}(a, g(x)) = x$, et $b = g(x)$ a pour image $d$.
\end{proof}

\begin{proof}[Démonstration des pages 99, 101 et 102]
Les applications $\varphi_\alpha$ sont bien définies, car $(\sigma x)_\alpha = \sigma_\alpha(x_\alpha)$ a même classe que $x_\alpha$.

(1) \emph{Injectivité.} Montrons d'abord que deux triangles supérieurs $x$, $y$ avec $y_\alpha = x_\alpha$ et $y_\beta \in \{x_\beta, \sigma_\beta x_\beta\}$ sont égaux. Si $y_\beta = \sigma_\beta x_\beta$, le point $x_\alpha = y_\alpha$ serait lié à $x_\beta$ et à $\sigma_\beta x_\beta$, contrairement à la proposition~2 ; donc $y_\beta = x_\beta$, et $y = x$ par l'unicité de~(4). Soient alors $x$, $y$ de mêmes images dans $E_\alpha$ et $E_\beta$. Si $y_\alpha = x_\alpha$, on conclut. Si $y_\alpha = \sigma_\alpha x_\alpha$, on remplace $y$ par $\sigma y$, qui a même classe : $(\sigma y)_\alpha = x_\alpha$ et $(\sigma y)_\beta$ est dans l'orbite de $x_\beta$, donc $\sigma y = x$. \emph{Surjectivité.} Soient $s \in \widetilde E_\alpha$ et $u \in \widetilde E_\beta$. Par la proposition~5, $u$ est lié à $s$ ou à $\sigma_\alpha s$, et~(4) donne un triangle supérieur dont les coordonnées d'indices $\alpha$ et $\beta$ sont ce point et $u$ ; sa classe a pour images les classes de $s$ et de $u$.

(2) Si $\bar x = \bar y$ et $x_\alpha = y_\alpha$, alors $y = x$ ou $y = \sigma x$, et le second cas donnerait $\sigma_\alpha x_\alpha = x_\alpha$. Si $(g, s)$ est dans le produit fibré, avec $g = \bar x$, la classe de $x_\alpha$ est celle de $s$, donc $s = x_\alpha$ ou $s = (\sigma x)_\alpha$ ; dans les deux cas $(g, s)$ est l'image d'un triangle supérieur.

(3) Par construction, $\psi_{\beta\alpha}$ envoie l'image de $x$ dans le premier produit fibré sur l'image de $x$ dans le second ; les deux ont pour première composante $\bar x$, et les relations de transitivité se lisent sur cette description.

(4) Les deux premières égalités sont le lemme~\ref{lem:74-revetement} pour $\sigma_\alpha$ et pour $\sigma$, la troisième est~(1). Pour la dernière, on applique le lemme~\ref{lem:74-ligne} à $f = \varphi_\alpha$, $g = \varphi_\beta$, $h = \varphi_\gamma$ et à la classe d'un point de $\widetilde E_\alpha$.
\end{proof}

\begin{proof}[Démonstration de la Scholie et de sa réciproque]
(1) L'application qui envoie $\alpha \in t_0$ sur $t_0(\alpha)$ et $s \in \widetilde E_\alpha$ sur $s$ est bijective par~(1) des pages~92 à~99. Elle respecte les arêtes : deux sommets de $t_0$ sont liés ; $t_0(\alpha)$ est lié à $s \in \widetilde E_\beta$ si et seulement si $\beta = \alpha$, par~(1) ; deux points d'un même $\widetilde E_\alpha$ sont liés si et seulement si l'un est l'image de l'autre par $\sigma_\alpha$, par~(2) ; deux points de $\widetilde E_\alpha$ et de $\widetilde E_\beta$, $\alpha \neq \beta$, sont liés dans $S$ si et seulement s'ils sont deux coordonnées d'un triangle supérieur, par~(4). Les conditions~(*) et $\beta_{\widetilde\Gamma}$ sont~(4) et~(3). Pour~(Ka), soient $x$, $y$, $z$ trois triangles supérieurs avec $x_\alpha = y_\alpha$, $x_\beta = z_\beta$, $y_\gamma = z_\gamma$. Le point $y_\gamma$ est lié à $y_\alpha = x_\alpha$, et, comme $z_\gamma$, à $z_\beta = x_\beta$. L'axiome~1 pour l'arête $\{x_\alpha, x_\beta\}$ donne $x_\gamma = y_\gamma$, et~(*) donne $x = y$.

(2) \emph{Axiome~1.} On examine les arêtes une à une. L'arête $\{t_0(\alpha), t_0(\beta)\}$ a pour seul troisième sommet $t_0(\gamma)$ : un point de $\widetilde E_\delta$ n'est lié qu'au sommet $t_0(\delta)$. L'arête $\{t_0(\alpha), s\}$ a pour seul troisième sommet $\sigma_\alpha s$ : un sommet lié à $t_0(\alpha)$ et à $s$ est dans $\widetilde E_\alpha$, et le seul point de $\widetilde E_\alpha$ lié à $s$ est $\sigma_\alpha s$. L'arête $\{s, \sigma_\alpha s\}$ a pour seul troisième sommet $t_0(\alpha)$ : un point $r$ lié à $s$ et à $\sigma_\alpha s$ qui ne serait pas dans $t_0$ ne pourrait être dans $\widetilde E_\alpha$, où il serait égal à $\sigma_\alpha s$ et à $s$, ni dans un autre $\widetilde E_\beta$, où $\beta_{\widetilde\Gamma}$ l'interdit. Enfin, soient $s \in \widetilde E_\alpha$ et $u \in \widetilde E_\beta$ liés par $\tau \in \widetilde T$. Le point $\tau_\gamma$ est lié à $s$ et à $u$. Soit $r$ un autre sommet lié à $s$ et à $u$. Il n'est pas dans $t_0$, car il serait lié à $t_0(\alpha)$ et à $t_0(\beta)$. S'il était dans $\widetilde E_\alpha$, il serait égal à $\sigma_\alpha s$, et $u$ serait lié à $s$ et à $\sigma_\alpha s$, contre $\beta_{\widetilde\Gamma}$ ; de même il n'est pas dans $\widetilde E_\beta$. Il est donc dans $\widetilde E_\gamma$, lié à $s$ par un $\rho_1 \in \widetilde T$ et à $u$ par un $\rho_2 \in \widetilde T$. Alors $\tau_\alpha = (\rho_1)_\alpha = s$, $\tau_\beta = (\rho_2)_\beta = u$ et $(\rho_1)_\gamma = (\rho_2)_\gamma = r$, et~(Ka) donne $\tau = \rho_1$, d'où $r = \tau_\gamma$.

Ce dernier argument montre aussi que trois points deux à deux liés hors de $t_0$ sont les coordonnées d'un même $\tau$ : deux d'entre eux ne sont pas dans une même orbite, car un point lié à $s$ et à $\sigma_\alpha s$ est dans $t_0$, et le troisième est le troisième sommet de l'arête des deux premiers.

\emph{Axiome~2.} Un triangle du graphe est de l'un des trois types suivants : $t_0$ ; un triangle vertical $\{t_0(\alpha), s, \sigma_\alpha s\}$ ; un triangle supérieur $\{\tau_\alpha\}$, $\tau \in \widetilde T$. Un triangle qui contient deux sommets de $t_0$ et un point de $\widetilde E_\alpha$ demanderait que ce point soit dans deux $\widetilde E$ ; un triangle formé de $t_0(\alpha)$ et de deux points est vertical, puisque ces points sont dans $\widetilde E_\alpha$ ; un triangle formé de trois points est supérieur par ce qui précède. Pour $t_0$, tout point de $\widetilde E_\alpha$ est lié à $t_0(\alpha)$. Pour un triangle vertical, un sommet de $t_0$ autre que $t_0(\alpha)$ est lié à $t_0(\alpha)$, ainsi qu'un point de $\widetilde E_\alpha$ ; un point $w \in \widetilde E_\beta$, $\beta \neq \alpha$, est lié à $s$ ou à $\sigma_\alpha s$ par $\beta_{\widetilde\Gamma}$. Pour un triangle supérieur $\tau$, un sommet $t_0(\delta)$ est lié à $\tau_\delta$. Soit $w \in \widetilde E_\delta$. Si $w = \sigma_\delta \tau_\delta$, il est lié à $\tau_\delta$. Sinon, notons $\beta$, $\gamma$ les deux autres indices, et supposons $w$ lié ni à $\tau_\beta$ ni à $\tau_\gamma$. Par $\beta_{\widetilde\Gamma}$, $\sigma_\delta w$ est lié à $\tau_\beta$ par un $\rho_1$ et à $\tau_\gamma$ par un $\rho_2$. La condition~(Ka) pour les indices $(\beta, \gamma, \delta)$ et les triplets $\tau$, $\rho_1$, $\rho_2$ donne $\tau = \rho_1$, d'où $\tau_\delta = \sigma_\delta w$ et $w = \sigma_\delta \tau_\delta$, contre l'hypothèse. C'est l'argument de la proposition~6, où la page invoque la proposition~2.

(3) Le sens réciproque est~(2). Supposons l'axiome~1, et soient $\tau$, $\rho_1$, $\rho_2$ comme dans~(Ka). Les points $\tau_\gamma$ et $(\rho_1)_\gamma$ sont liés à $\tau_\alpha$ (par $\tau$ et par $\rho_1$) et à $\tau_\beta$ (par $\tau$ et par $\rho_2$). L'axiome~1 pour l'arête $\{\tau_\alpha, \tau_\beta\}$ donne $\tau_\gamma = (\rho_1)_\gamma$, et~(*) pour les indices $\alpha$, $\gamma$ donne $\tau = \rho_1$.

(4) Si $\tau$ et $\rho$ ont mêmes coordonnées d'indices $\alpha$ et $\beta$, (Ka) pour les indices $(\alpha, \gamma, \beta)$ et les triplets $\tau$, $\rho$, $\tau$ donne $\tau = \rho$. (Ka) entraîne (Ka$'$) : pour $(s_\alpha)$ et $(t_\alpha)$ comme dans (Ka$'$), (Ka) appliquée à $t_c$, $t_b$, $t_a$ donne $t_c = t_b$, et $(s_\alpha)$ a les coordonnées de $t_c$. Inversement, sous~(*) et (Ka$'$), soient $\tau$, $\rho_1$, $\rho_2$ comme dans (Ka). La famille $s$ égale à $\tau$ hors de $\gamma$ et à $(\rho_1)_\gamma$ en $\gamma$, avec $t_\gamma = \tau$, $t_\beta = \rho_1$, $t_\alpha = \rho_2$, vérifie les hypothèses de (Ka$'$) ; donc $s \in \widetilde T$, puis $s = \tau$ par~(*), d'où $\tau_\gamma = (\rho_1)_\gamma$ et $\tau = \rho_1$ par~(*).
\end{proof}

\begin{proof}[Démonstration des problèmes I et II]
\emph{(*).} Si $\psi_\alpha(x) = \psi_\alpha(y)$ et $\psi_\beta(x) = \psi_\beta(y)$, alors $\varphi_\alpha \pi(x) = p_\alpha \psi_\alpha(x) = \varphi_\alpha \pi(y)$, et de même pour $\beta$ ; $\beta_\Gamma$ donne $\pi(x) = \pi(y)$, et l'injectivité de $(\pi, \psi_\alpha)$ donne $x = y$.

\emph{$\beta_{\widetilde\Gamma}$.} Soient $s \in \widetilde E_\alpha$ et $u \in \widetilde E_\beta$. Deux éléments $x$, $y$ de $\widetilde\Gamma$ avec $\psi_\beta(x) = \psi_\beta(y) = u$ et $p_\alpha\psi_\alpha(x) = p_\alpha\psi_\alpha(y) = p_\alpha(s)$ ont la même image dans $E_\alpha \times E_\beta$, donc la même image dans $\Gamma$ par $\beta_\Gamma$, et sont égaux par l'injectivité de $(\pi, \psi_\beta)$. Si $s$ et $\sigma_\alpha s$ étaient tous deux liés à $u$, par $x$ et par $y$, on aurait donc $x = y$ et $s = \sigma_\alpha s$. Pour l'existence, soit $g \in \Gamma$ d'image $(p_\alpha(s), p_\beta(u))$, et $x$ au-dessus de $g$ avec $\psi_\beta(x) = u$ ; alors $p_\alpha \psi_\alpha(x) = \varphi_\alpha(g) = p_\alpha(s)$, donc $\psi_\alpha(x) \in \{s, \sigma_\alpha s\}$.

La suite résulte de la réciproque de la Scholie. Pour un graphe, les données de l'énoncé sont $\Gamma$, les $E_\alpha$, les classes $\widetilde E_\alpha \to E_\alpha$, $\widetilde\Gamma$, $x \mapsto \bar x$ et $x \mapsto x_\alpha$ ; les hypothèses sont~(1) et~(2) de l'énoncé des pages~99, 101 et~102.
\end{proof}

\begin{proof}[Démonstration de la nécessité de (Ka)]
Les conditions~(*) et $\beta_{\widetilde\Gamma}$ portent sur un nombre fini de cas et sont vérifiées par énumération. Les triplets $\tau = ((0,0), (0,0), (0,0))$, $\rho_1 = ((0,0), (1,0), (1,0))$ et $\rho_2 = ((1,0), (0,0), (1,0))$ de $\widetilde T$ vérifient $\tau_a = (\rho_1)_a$, $\tau_b = (\rho_2)_b$ et $(\rho_1)_c = (\rho_2)_c$, et $\tau \neq \rho_1$. La condition~(Ka) est donc fausse, et l'axiome~1 aussi par~(3) de la Scholie.
\end{proof}

\begin{proof}[Démonstration de la page 116]
Si les $\varphi_a$ préservent la relation, on l'applique à $(1, 1, 1)$ et à $a = g^{-1}$ : il vient $g \cdot g = 1$. Inversement, si $g^2 = 1$ pour tout $g$, alors $g^{-1} = g$, et $gk = (gk)^{-1} = k^{-1}g^{-1} = kg$ : le groupe est commutatif. Pour $xyz = 1$, on a $x \cdot a^{-1}z^{-1} \cdot a^{-1}y^{-1} = x \cdot az \cdot ay = a^2 \cdot xyz = 1$.
\end{proof}

\begin{remarque}
La direction directe n'emploie l'axiome~2 que pour $t_0$ et pour les triangles verticaux. Elle n'emploie pas l'axiome~3 de la page~100 ; celui-ci ne sert qu'aux dénombrements. Les cas exclus par l'axiome~3, un triangle seul ou une étoile de triangles en un sommet, ont $\widetilde\Gamma = \varnothing$, et leur $\Gamma$ vide vérifie $\beta_\Gamma$ puisque deux des $E_\alpha$ au moins sont vides. La page~101, qui demande $\widetilde\Gamma \neq \varnothing$, les écarte.
\end{remarque}

\begin{remarque}
La page~97 affirme que des données vérifiant~(*) donnent un graphe qui vérifie l'axiome~1. Le contre-exemple de la nécessité de (Ka) vérifie~(*), (**) et $\beta_{\widetilde\Gamma}$, et son graphe ne vérifie pas l'axiome~1. L'énoncé de la page~97 demande donc la condition de la marge, que l'accolade rattache à~(*) et~(**). Sous~(*), l'axiome~1 du graphe des données équivaut à (Ka), et cette équivalence permet de restituer la condition. La conclusion de la marge est illisible : les formes (Ka) et (Ka$'$) sont les deux lectures, équivalentes sous~(*), qui rendent la Scholie juste. La condition~(*) n'est pas utilisée pour la réciproque ; elle découle de~(Ka).
\end{remarque}

\begin{remarque}
Pour l'axiome~2 aux triangles supérieurs, la page~100 (proposition~6) invoque la proposition~2, qui vient de l'axiome~1 dans un graphe donné. Dans la réciproque, la démonstration utilise (Ka) à cet endroit.
\end{remarque}

\begin{remarque}
L'énoncé de la page~116 suppose dès le départ que le carré latin est la table d'un groupe. La lecture écrit qu'il explique pourquoi le carré latin du problème~I \q{doit être la table d'un $\mathbb F_2$-espace vectoriel}. L'énoncé démontré ici est plus étroit : si $\Gamma$ est la table de $xyz = 1$ dans un groupe et si les $\varphi_a$ la préservent, le groupe est un $2$-groupe abélien élémentaire. La page applique la relation à $y = z = 1$ ; la démonstration prend $x = y = z = 1$, ce qui revient au même puisque $x$ est déterminé par $y$ et $z$.
\end{remarque}

Ne sont pas démontrés ici : la classification des graphes vérifiant les axiomes~1 à~3 et leur identification aux quadrangles généralisés d'ordre $(2, c)$, la description des carrés latins du problème~I et la classification des revêtements doubles du problème~II, leur lien avec les extensions centrales des pages~32 à~39, l'affirmation de la lecture que le carré latin est en général la table d'un $\mathbb F_2$-espace vectoriel, et les nombres de la page~104 pour $|S|$, $|A|$ et $|T|$. L'identification des $\widetilde E_\alpha$ du graphe des données aux ensembles de départ, qui achèverait l'équivalence dans le sens réciproque, n'est démontrée que pour les triangles supérieurs.
